\section{Further research: precise questions and next steps}
\label{sec:research}

The results above close several concrete gaps, but each also
identifies a natural boundary. The following programme separates
conjectures with specific existing evidence from broader questions.
The proposed tasks are intended to lead to additional theorems
or falsifiable identities, not merely more decimal tables.

\subsection{Euler bounds: rates, extremizers, and further moment families}

\paragraph{Sharp bounds for a fixed truncation index.}
The global constant $C_*$ is forced by $N=1$ and $a\downarrow0$.
For practical computations with large $N$, it leaves substantial
room. Define the exact pointwise envelope
\[
 K_N(v)=\max_{0\leq x\leq1}\{W_N(vx)-W_N(x)\},
 \qquad 0<v<1.
\]
Theorem~\ref{thm:kernel} determines the supremum over all $N$,
but does not determine $K_N(v)$ in a useful closed form for
$N\geq2$. Find its maximizing point, its dependence on $N$,
and the strongest integrated bound it yields for fixed $a,b$.
A derivative equation for the maximizing point is elementary,
so the challenge is an interpretable exact inequality rather
than a numerical maximization.

The scaling $x=y/N$ suggests the limiting comparison
$e^{-vy}-e^{-y}$, whose maximum is
$(1-v)v^{v/(1-v)}$. This suggests asking for an error-controlled
limit of $K_N(v)$, including its behavior as $v\downarrow0$
and $v\uparrow1$. Uniformity at those endpoints matters
before integrating against an order-dependent measure.
The elementary limiting calculation alone does not supply
that uniformity.

\paragraph{Fractional-order late Euler errors.}
The integer-order manuscript already gives a polynomial in
$\tfrac12\log N$ multiplying $2^{-N}N^{-1/2}$, with leading
coefficient determined by the total weight. Extend the full
expansion to positive nonintegral $a,b$, explicitly tracking
the extra endpoint terms that appear when an exponent crosses
an integer. In particular, establish a uniform version of the
fixed-order leading behavior suggested by the two-measure
convolution,
\[
 E_N-g_{a,b}\sim
 \frac{\sqrt{\pi}}{2^{a+b}\Gamma(a+b)}
 \frac{2^{-N}(\log N)^{a+b-1}}{\sqrt N}.
\]
For positive integral orders this is already a theorem in
\src{chapters/05-signed-kernels.tex}; the research target is
the fractional continuation and its uniform error structure.
A proof should separate the region where both logarithmic
moment variables are large from the boundary layers in which
only one is large. The signed-measure representation can fail
to have finite mass, so the positive two-measure difference
is a promising starting point.

\paragraph{The rate of approach to the asymptotic global constant.}
Theorem~\ref{thm:euler-asymptotic-budget} proves $C_N\to1$, and
Proposition~\ref{prop:euler-axis-lower} gives a concrete lower
asymptotic. The natural next target is the matching upper bound.

\begin{conjecture}[Asymptotic optimal Euler constant]
\label{conj:euler-rate}
With
\[
 p=\frac{\log2}{\log(3/2)},\qquad
 q=\frac{\log3}{\log2},\qquad c=(1-q^{-1})q^{-p},
\]
the optimal scaled-error constant satisfies
\[
 C_N-1\sim cN^{-p}.
\]
\end{conjecture}

The matching lower bound is proved. The global upper bound over
positive $a$ is open, and even the sharp upper asymptotic for
the supremum over all $b$ on the axis requires an argument.
The moving-order phase law supplies a broad region with $R_N\to1$,
but that limit does not distinguish errors of order $N^{-p}$.
Any statement about near maximizers must specify the accuracy:
an error $o(C_N-1)$ is relevant here, whereas an error merely
$o(1)$ need not force $a\to0$. Monotonicity of $C_N$ is also
not asserted.

The diagnostic program \src{code/axis_remainder_diagnostics.py}
evaluates finite Euler sums at $210$ decimal digits.
For the following indices it finds numerical axis stationary
points, with the displayed scaled excess:
\begin{center}
\small
\begin{tabular}{@{}rrrr@{}}
\toprule
$N$ & $\widehat b_N$ & Stationary $b$ &
$(R_N(0,b)-1)/(cN^{-p})$\\
\midrule
4   & 4.5549051506  & 4.7024351462  & 0.8809410\\
16  & 7.9739277333  & 7.8500957399  & 1.0384240\\
64  & 11.3929503160 & 11.2260085539 & 1.0760668\\
256 & 14.8119728987 & 14.6711951089 & 1.0684372\\
\bottomrule
\end{tabular}
\end{center}
Their global optimality is not certified. At $b=\widehat b_N$,
the three tested positive values $a=10^{-4},10^{-2},10^{-1}$
all give smaller remainders than $a=0$. These are limited axis
and local checks; they do not establish the conjecture or
monotone convergence of the displayed ratios. An effective
compact-rectangle estimate, combined with a sharp analysis
of the axis scaling, is one possible route to an upper bound.

\paragraph{Further moving-order boundaries.}
Theorem~\ref{thm:euler-moving} treats two orders growing
proportionally to $\log N$. Determine the finer behavior
when one of the leading coefficients $A,B$ tends to zero,
or when one order remains fixed while the other crosses
the logarithmic threshold. The bounded-width gamma variable
then interacts with the smoothed threshold itself, and the
two normal profiles need not describe the first correction.
Also study orders tending to zero, where endpoint
concentration affects the fractional late-error expansion.
A useful result would connect the several scaling regimes
with explicit overlapping error bounds.

\paragraph{Other moment products.}
Theorem~\ref{thm:moment-general} applies to any two finite
positive measures. Explore shifted harmonic numbers,
Hurwitz-type inner sums, mixtures of orders, and products
arising from completely monotone sequences. For each family,
identify the sharp measure-specific constant rather than
automatically using $(1+\sqrt2)/2$.
Higher-depth nested partial sums lead to higher divided
differences of the same kernel. Determine which of those
differences have a fixed sign and what sharp moduli replace
the two-point inequality. This is a concrete route toward
certified Euler algorithms for real-order higher-depth
polylogarithms.

\subsection{Formal depth and genuine functional reduction}

\paragraph{Exploit the full binary formula constructively.}
The finite-direction dimension theorem gives the rank but
does not choose the best basis for an identity search.
Construct rational bases adapted to the irreducible
$\operatorname{Sym}^{w-2j}V\otimes\det^j$ summands and
their confluent interpolation conditions. This should make
it possible to solve reduction problems without forming
all $2^w$ words. The weight-seven product correction is
explicit but lengthy; a basis respecting reversal,
complementation, and endpoint regularization may produce a
much shorter analytic expression.

\paragraph{Compare with the Nielsen literature at each weight.}
The height-one results must be compared with the known
Nielsen reductions, including their meaning modulo products
and their exact symbol spaces. A useful next deliverable is
a table that gives, for every weight in a practical range,
the published reduction, a word certificate in the repository's
conventions, and the endpoint correction on a stated branch
domain. It should identify equivalence of relations, not
count the same identity multiple times because of a change
of notation.

\paragraph{Go beyond three punctures.}
The six transformations act on only three projective letter
directions. A general rational substitution may introduce
additional singularities and therefore a larger alphabet;
it is not represented by an arbitrary $2\times2$ change of
the old letters. Extend the rank problem to a specified
puncture set and a specified collection of rational maps.
Describe the kernel of specialization back to the original
periods. The weight-eight obstruction here is a useful test
case: which genuinely new pullbacks or distribution
relations remove it, and what new constants do they introduce?

\paragraph{Separate formal and period questions.}
For the unresolved Gaussian $S_6$ and $S_8$ evaluations,
specify the expected algebra of constants before searching.
Construct the corresponding functional relation space,
map it to the target specialization, and compare the proposed
identity with that image. A formal failure would indicate
that the chosen tools are insufficient. It would not prove
period independence. Conversely, an exact functional
certificate with controlled endpoint constants would settle
an evaluation without relying on an integer-relation residual.

\subsection{The global radial phase portrait}

Let $a_*(b)$ be the already proved unique solution of $K=0$
in $(1,2)$ for $0<b<1$, and put $q(b)=Q(a_*(b),b)$.
The pinned manuscript proves opposite signs at two rational
parameters and reports a wider diagnostic scan. The exact
certificate here proves one simple zero near the upper
endpoint of that interval.

\begin{conjecture}[Unique quartic transition on the quadratic threshold]
\label{conj:quartic-global}
The function $q$ has exactly one zero on $0<b<1$, namely
the certified $b_\dagger$. It is negative for
$0<b<b_\dagger$ and positive for $b_\dagger<b<1$.
\end{conjecture}

The local sign at the certified zero is proved, and the
existing finite scan supports the global statement. Neither
fact excludes additional zeros away from the isolating box.
A rigorous approach should combine asymptotic analysis as
$b\downarrow0$ and $b\uparrow1$ with a finite interval cover
of the remaining compact region. Direct interval evaluation
may be inefficient near $b=1$, where several low-order
coefficients nearly cancel. Factoring the known degeneration
at $(a,b)=(1,1)$ before bounding remainders is likely to be
essential.

\paragraph{Global continuation of the two-extremum region.}
The local fold and the explicit rational example prove that
a minimum followed by a maximum can occur. Follow the two
critical branches to larger radii and determine how they
end: collision, escape through $\rho=1$, or interaction with
another branch. Can more than two radial critical points
occur for some positive orders? The local concavity proof
does not address this global question. A plausible first
step is an interval continuation of the coupled equations
for the angular zero and its radial derivative, with a
separate proof that no branch is missed.

\paragraph{The older integer-order monotonicity question.}
The conjectured full-radius decrease for integer $a\geq2$,
and the special $a=1$ behavior, remain interesting.
The fractional examples do not refute them: their parameters
are close to, but different from, $(1,1)$.
Look for a sign-definite integral expression for the radial
derivative in the integer range. A proof based only on the
coefficient $K$ would establish local behavior and leave the
rest of the radius interval untreated.

\paragraph{Higher degeneracies.}
The binomial implicit equation gives every radial coefficient
without guessing higher Chebyshev expansions. Search for
families with $K=Q=R=0$ after adding a deformation parameter,
such as a Lerch shift or an extra inner moment. If a first
nonzero coefficient appears at $\rho^{2\ell}$, generic
one-parameter perturbation suggests turning radii with
power $1/(2\ell-2)$. State the transversality hypotheses
explicitly and prove the local normal form before drawing
a global bifurcation diagram.

\subsection{Growing reflected exponents and optimal truncation}

\paragraph{The maximal proportional-exponent regime.}
The shifted-saddle continuation establishes a nonzero
interval of proportional growth. Determine the largest
admissible interval by analyzing the real saddle equation
and competing complex saddles. Positivity of the Taylor
coefficients of $\Gamma(1-u)$ is available, but
$-\log\Gamma(1+u)$ does not have a positive coefficient
sequence. A local stationary point alone therefore does
not prove global contour dominance for all $m/k$.
Locate the first loss of dominance or the first saddle
collision, and give an exact or certified description of
the transition.

\paragraph{Move the optimal truncation index as well.}
When $m/N$ is no longer small, the exponential rate of the
coefficients changes. The optimal truncation index must
then be determined from that joint rate, rather than
retaining $N\Lambda\approx n$ without adjustment.
Derive the correct stationarity equation when $n,m,N$
grow together, and test whether a half-term law persists.
The signed jump on the negative cut requires absolute
moment estimates; eventual coefficient signs do not by
themselves make the cut density positive.

\paragraph{Effective constants and preasymptotic behavior.}
Turn the present uniform big-$O$ statements into explicit
inequalities in a chosen practical range. The diagnostic
examples show why this matters: a fixed range parameter
can be mathematically admissible while its asymptotic
behavior appears only at much larger indices.
Effective contour radii, bounds on logarithmic derivatives,
and a computable separation of the central arc from the
outer contour would turn the theorem into a usable
certified truncation prescription.

\paragraph{Connect the different moment scales.}
The coefficient scale $m\asymp\sqrt{k}$ and the direct-moment
scale involving $n/(m+1)-\log m$ arise from different
representations. Determine their overlap after the
optimal index is chosen, and compare the resulting
asymptotic expansions term by term. A matching theorem
would explain which representation is stable in each
region and where cancellations make one numerically
preferable.

\subsection{A practical order of attack}

The most bounded next tasks are an exact global sign study
for $q(b)$, a basis-efficient implementation of the finite-direction
rank theorem, and effective constants for the shifted saddle.
Each has a clearly specified output and can be checked
against the certificates in this article. The longer
projects are a global radial classification, an extension
to larger puncture alphabets, and a joint optimal-truncation
theory for proportional reflected exponents. The special-value
conjectures should remain visible throughout, but their
status should change only when an exact functional or
arithmetic certificate becomes available.
